\chapter{ปัญญาประดิษฐ์และการเรียนรู้เชิงลึกสำหรับเกษตรกรรมแม่นยำ}
\label{chap:deeplearning}

หัวใจสำคัญของการยกระดับระบบสมาร์ตฟาร์มจากการเป็นเพียง \textbf{"ระบบเฝ้าระวัง (Monitoring)"} ไปสู่การเป็น \textbf{"ระบบพยากรณ์และตัดสินใจอัจฉริยะ (Predictive Autonomy)"} คือการประยุกต์ใช้โมเดลการเรียนรู้เชิงลึก (Deep Learning)\index{Deep Learning} เพื่อวิเคราะห์ความสัมพันธ์แบบไม่เป็นเชิงเส้น (Non-linear Dynamics) ระหว่างสภาวะบรรยากาศ แสงแดด และการสูญเสียน้ำในดิน บทนี้นำเสนอทฤษฎีและการประยุกต์ใช้โครงข่ายประสาทเทียมแบบหน่วยความจำระยะสั้นยาว (LSTM)\index{LSTM} ร่วมกับชุดข้อมูลที่เก็บได้จริงจากระบบ JC -AgriTech + AI

\section{ทฤษฎีการพยากรณ์อนุกรมเวลาหลายตัวแปรในแปลงเกษตร}

พลวัตความชื้นในดินบริเวณเขตรากพืชขึ้นอยู่กับปฏิสัมพันธ์อันซับซ้อนระหว่างกระบวนการระเหยของน้ำจากผิวดิน (Evaporation) และการคายน้ำของปากใบพืช (Transpiration) ซึ่งรวมเรียกว่า \textbf{การคายระเหยน้ำ (Evapotranspiration หรือ ET)}\index{Evapotranspiration (ET)} ตามสมการมาตรฐาน FAO-56 Penman-Monteith\index{Penman-Monteith}

\begin{equation}
    ET_0 = \frac{0.408 \Delta (R_n - G) + \gamma \frac{900}{T+273} u_2 (e_s - e_a)}{\Delta + \gamma(1 + 0.34 u_2)}
    \label{eq:penman_monteith}
\end{equation}

โดยที่:
\begin{itemize}
    \item $R_n$ คือรังสีสุทธิที่ผิวพืช (Net Radiation, $\text{MJ}\cdot\text{m}^{-2}\cdot\text{day}^{-1}$) ซึ่งวัดได้จากเซนเซอร์โดมตะวัน BH1750
    \item $T$ คืออุณหภูมิอากาศเฉลี่ย ($^\circ\text{C}$) วัดจากเซนเซอร์ SHT45
    \item $(e_s - e_a)$ คือค่าแรงดึงระเหยน้ำของอากาศ (Vapor Pressure Deficit: VPD, $\text{kPa}$)
    \item $\Delta$ คือความชันของเส้นโค้งความดันไออิ่มตัว ($\text{kPa}\cdot{^\circ\text{C}}^{-1}$)
    \item $\gamma$ คือค่าคงที่ไซโครเมตริก ($\text{kPa}\cdot{^\circ\text{C}}^{-1}$)
\end{itemize}

อัตราการเปลี่ยนแปลงความชื้นในดินบริเวณเขตรากพืช ($\theta$) สามารถจำลองด้วยสมการสมดุลน้ำ (Water Balance Equation):

\begin{equation}
    \frac{d\theta}{dt} = \frac{1}{Z_r} \Big( I(t) - ET(t) - D(t) \Big)
    \label{eq:water_balance}
\end{equation}

เมื่อ $Z_r$ คือความลึกเขตรากพืช, $I(t)$ คืออัตราการให้น้ำจากปั๊มรีเลย์, และ $D(t)$ คือการซึมลึก โจทย์การพยากรณ์จึงเป็นการทำนายค่า $\hat{y}_{t+H}$ (ความชื้นในดินล่วงหน้า $H$ ก้าวเวลา) จากประวัติข้อมูลย้อนหลัง $W$ ก้าวเวลา

\begin{equation}
    \hat{y}_{t+H} = f_\theta\Big( \mathbf{x}_{t-W+1}, \mathbf{x}_{t-W+2}, \dots, \mathbf{x}_t \Big)
    \label{eq:multivariate_forecasting}
\end{equation}

โดยที่ $\mathbf{x}_t \in \mathbb{R}^D$ คือเวกเตอร์คุณลักษณะขนาด $D=12$ มิติ ณ เวลา $t$

\section{สถาปัตยกรรมโครงข่ายประสาทเทียมแบบ LSTM}

โครงข่ายประสาทเทียมแบบ \textbf{LSTM (Long Short-Term Memory)} ได้รับการออกแบบให้มีเซลล์หน่วยความจำ (Memory Cell $c_t$) และกลไกประตูกำกับสัญญาณ 3 ประตู ได้แก่ ประตูลืม (Forget Gate $f_t$) ประตูรับข้อมูล (Input Gate $i_t$) และประตูส่งออก (Output Gate $o_t$) ดังแสดงในภาพที่ \ref{fig:tikz_lstm_unrolled}

\begin{figure}[htbp]
\centering
\resizebox{0.86\textwidth}{!}{
\begin{tikzpicture}[
    node distance=1.2cm and 1.2cm,
    lstm/.style={rectangle, draw=rbruNavy, fill=agriBlue!10, thick, rounded corners=4pt, align=center, minimum height=1.5cm, minimum width=2.2cm, drop shadow={opacity=0.15}},
    inputnode/.style={circle, draw=agriGreen, fill=agriGreen!15, thick, inner sep=2pt, minimum size=0.9cm, font=\small},
    outputnode/.style={circle, draw=agriRed, fill=agriRed!15, thick, inner sep=2pt, minimum size=0.9cm, font=\small},
    dense/.style={rectangle, draw=rbruGold!90!black, fill=rbruGold!20, thick, rounded corners=3pt, align=center, minimum height=0.9cm, minimum width=2.2cm},
    line/.style={draw=rbruNavy, -latex', thick},
    stateline/.style={draw=agriBlue!80!black, -latex', thick}
]
    % Input Nodes
    \node[inputnode] (x1) {$\mathbf{x}_{t-2}$};
    \node[inputnode, right=2.8cm of x1] (x2) {$\mathbf{x}_{t-1}$};
    \node[inputnode, right=2.8cm of x2] (x3) {$\mathbf{x}_t$};

    % LSTM Cell Blocks
    \node[lstm, above=1.0cm of x1] (lstm1) {\textbf{LSTM Cell}\\\footnotesize $\mathbf{c}_{t-2}, \mathbf{h}_{t-2}$};
    \node[lstm, above=1.0cm of x2] (lstm2) {\textbf{LSTM Cell}\\\footnotesize $\mathbf{c}_{t-1}, \mathbf{h}_{t-1}$};
    \node[lstm, above=1.0cm of x3] (lstm3) {\textbf{LSTM Cell}\\\footnotesize $\mathbf{c}_t, \mathbf{h}_t$};

    % Dense Regressor
    \node[dense, above=1.0cm of lstm3] (dense_reg) {\textbf{Dense Regressor}\\\footnotesize Linear(64 $\to$ 1)};
    \node[outputnode, above=0.8cm of dense_reg] (y_pred) {$\hat{y}_{t+H}$};

    % Connections Inputs to LSTM
    \draw[line] (x1) -- (lstm1);
    \draw[line] (x2) -- (lstm2);
    \draw[line] (x3) -- (lstm3);

    % Recurrent Connections
    \draw[stateline] (lstm1) -- node[above, font=\scriptsize] {$\mathbf{c}_{t-2}, \mathbf{h}_{t-2}$} (lstm2);
    \draw[stateline] (lstm2) -- node[above, font=\scriptsize] {$\mathbf{c}_{t-1}, \mathbf{h}_{t-1}$} (lstm3);

    % Output Pipeline
    \draw[line] (lstm3) -- node[right, font=\scriptsize] {$\mathbf{h}_t$ (State)} (dense_reg);
    \draw[line] (dense_reg) -- (y_pred);

\end{tikzpicture}
}
\caption{สถาปัตยกรรมโครงข่ายประสาทเทียมแบบคลี่ขยายตามเวลา (Unrolled LSTM Multi-step Architecture)}
\label{fig:tikz_lstm_unrolled}
\end{figure}

สมการการทำงานภายในเซลล์ LSTM กำหนดโดย:
\begin{align}
    f_t &= \sigma\Big( W_f [\mathbf{x}_t, h_{t-1}] + b_f \Big) \label{eq:forget_gate} \\
    i_t &= \sigma\Big( W_i [\mathbf{x}_t, h_{t-1}] + b_i \Big) \label{eq:input_gate} \\
    \tilde{c}_t &= \tanh\Big( W_c [\mathbf{x}_t, h_{t-1}] + b_c \Big) \label{eq:candidate_cell} \\
    c_t &= f_t \odot c_{t-1} + i_t \odot \tilde{c}_t \label{eq:cell_state} \\
    o_t &= \sigma\Big( W_o [\mathbf{x}_t, h_{t-1}] + b_o \Big) \label{eq:output_gate} \\
    h_t &= o_t \odot \tanh(c_t) \label{eq:hidden_state}
\end{align}

\section{การสร้างแบบจำลองและลูปการฝึกโครงข่ายด้วย PyTorch}

การฝึกโมเดลในไฟล์ \texttt{ai\_colab\_starter.py} เริ่มต้นจากการแปลงข้อมูลดิบเป็นหน้าต่างเลื่อนขนาด $W = 12$ สเต็ป (ข้อมูลย้อนหลัง 3 ชั่วโมง) เพื่อพยากรณ์ความชื้นในดินเขตรากพืชในอีก 1 ชั่วโมงข้างหน้า ดังแสดงในโค้ด PyTorch ต่อไปนี้

\begin{pythonbox}[การสร้างโครงข่ายประสาทเทียม AgriLSTMForecaster ด้วย PyTorch]
import torch
import torch.nn as nn
import torch.optim as optim

class AgriLSTMForecaster(nn.Module):
    """
    Multivariate LSTM Neural Network for Root-Zone Moisture Prediction.
    Designed by Asst. Prof. Dr. Chewa Thassana (JC-AgriTech + AI).
    """
    def __init__(self, input_dim=12, hidden_dim=64, num_layers=2):
        super().__init__()
        self.lstm = nn.LSTM(
            input_size=input_dim,
            hidden_size=hidden_dim,
            num_layers=num_layers,
            batch_first=True,
            dropout=0.2
        )
        self.fc = nn.Sequential(
            nn.Linear(hidden_dim, 32),
            nn.ReLU(),
            nn.Linear(32, 1) # Predict target root-zone moisture (%)
        )
        
    def forward(self, x):
        # x shape: [batch_size, seq_len, input_dim]
        out, (hn, cn) = self.lstm(x)
        prediction = self.fc(out[:, -1, :])
        return prediction

def train_agri_model(model, train_loader, val_loader, epochs=50, lr=1e-3):
    criterion = nn.MSELoss()
    optimizer = optim.Adam(model.parameters(), lr=lr, weight_decay=1e-5)
    
    for epoch in range(epochs):
        model.train()
        train_loss = 0.0
        for batch_x, batch_y in train_loader:
            optimizer.zero_grad()
            preds = model(batch_x)
            loss = criterion(preds, batch_y)
            loss.backward()
            torch.nn.utils.clip_grad_norm_(model.parameters(), max_norm=1.0)
            optimizer.step()
            train_loss += loss.item()
            
        # Validation step
        model.eval()
        val_loss = 0.0
        with torch.no_grad():
            for vx, vy in val_loader:
                vpreds = model(vx)
                val_loss += criterion(vpreds, vy).item()
                
        if (epoch + 1) % 10 == 0:
            print(f"Epoch [{epoch+1}/{epochs}] - Train Loss: {train_loss/len(train_loader):.4f} - Val Loss: {val_loss/len(val_loader):.4f}")
\end{pythonbox}

\section{สถาปัตยกรรม TinyML บนชิป ESP32-S3: การชดเชยค่า NPK, pH และความชื้นดิน (Soil Neural Calibrator)}
\label{sec:tinyml_soil_calibrator}

ในการใช้งานจริงภาคสนาม เซนเซอร์วัดคุณสมบัติดินแบบมัลติโพรบ 7-in-1 (RS485 Modbus RTU) และเซนเซอร์วัดความชื้นชนิดคาปาซิทีฟ (Soil Stick) มักเผชิญกับปัญหาข้อจำกัดทางฟิสิกส์และเคมีของดินอย่างรุนแรง ซึ่งสมการเชิงเส้นแบบดั้งเดิม ($y = mx + c$) ไม่สามารถแก้ไขได้ ระบบ JC -AGRITecH + AI จึงได้บูรณาการโครงข่ายประสาทเทียมขนาดเล็ก \textbf{(TinyML Deep Neural Calibrator)}\index{TinyML} ฝังลงในหน่วยความจำของชิปไมโครคอนโทรลเลอร์ ESP32-S3 เพื่อทำหน้าที่ประมวลผลชดเชยค่าความคลาดเคลื่อนแบบเรียลไทม์ (Edge AI Inference) ณ ปลายทาง

\subsection{ปัญหาทางฟิสิกส์-เคมีดินและการแทรกแซงข้ามมิติ (Cross-Sensitivity)}

ความคลาดเคลื่อนของเซนเซอร์ดินเกิดจาก 3 ปัจจัยทางธรรมชาติที่เกี่ยวพันกันแบบไม่เป็นเชิงเส้น:

\begin{enumerate}
    \item \textbf{การแทรกแซงของความชื้นต่อสภาพนำไฟฟ้า (EC-Moisture Coupling):} 
    ตามแบบจำลองไดอิเล็กทริกของ Hilhorst (2000)\index{Hilhorst Model} และสมการ Rhoades et al. (1976)\index{Rhoades Model} สภาพนำไฟฟ้ารวมในเนื้อดิน (Bulk EC: $\sigma_b$) จะแปรผันตามสัดส่วนปริมาตรน้ำในช่องว่างดิน ($\theta$) เมื่อดินแห้ง ฟิล์มน้ำรอบเม็ดดินจะขาดช่วง ส่งผลให้ความต้านทานไฟฟ้าพุ่งสูงและสภาพนำไฟฟ้าดิ่งลง เซนเซอร์วัด NPK จึงอ่านค่าธาตุอาหารตกวูบเป็นศูนย์ แม้ในดินจะมีเม็ดปุ๋ยสะสมอยู่เต็มก็ตาม
    
    \item \textbf{ผลกระทบจากอุณหภูมิดินต่อการเคลื่อนที่ของไอออน (Thermal Drift):}
    ตามสมการ Stokes-Einstein เมื่อดินมีอุณหภูมิสูงขึ้น ความหนืดของน้ำในช่องว่างดินจะลดลง ส่งผลให้ไอออนของธาตุอาหาร ($\text{NH}_4^+$, $\text{NO}_3^-$, $\text{H}_2\text{PO}_4^-$, $\text{K}^+$) เคลื่อนที่ได้คล่องตัวขึ้น (Ion Mobility สูงขึ้น) โดยตามมาตรฐาน USDA Handbook 60 ค่าสภาพนำไฟฟ้าจะขยับขึ้นประมาณ \textbf{$+1.9\%$ ถึง $+2.1\%$ ต่อ $1^\circ\text{C}$}:
    \begin{equation}
        \text{EC}_{T} = \text{EC}_{25} \times \Big[ 1 + \alpha(T_{\text{soil}} - 25) \Big]
        \label{eq:thermal_drift}
    \end{equation}
    หากในแปลงปลูกมีอุณหภูมิดินต่างกันระหว่างเช้าตรู่ ($20^\circ\text{C}$) และบ่ายแดดจัด ($38^\circ\text{C}$) ค่า NPK ดิบจะแกว่งตัวสูงเกินจริงถึง $+36\%$
    
    \item \textbf{การเลื่อนของความชันศักย์ไฟฟ้ากรด-ด่าง (Nernstian pH Drift):}
    ตามสมการเคมีไฟฟ้า Nernst ค่าความชันในการวัดกรด-ด่าง ($S = \frac{2.303 RT}{F}$) จะเพิ่มขึ้นจาก $59.16\,\text{mV/pH}$ ที่ $25^\circ\text{C}$ เป็น $62.13\,\text{mV/pH}$ ที่ $40^\circ\text{C}$ ทำให้การคำนวณค่า pH ดิบเบี่ยงเบนไปตามอุณหภูมิของแปลงปลูก
\end{enumerate}

\subsection{สถาปัตยกรรมโครงข่ายประสาทเทียม Multi-Layer Perceptron (MLP)}

เพื่อปลดล็อกการแทรกแซงข้ามมิติดังกล่าว ระบบได้ออกแบบโครงข่ายประสาทเทียมแบบ Multi-Layer Perceptron (MLP)\index{Multi-Layer Perceptron} ขนาดกะทัดรัด (10 Inputs $\to$ Dense 16 $\to$ Dense 16 $\to$ 5 Outputs) ดังโครงสร้าง:

\begin{equation}
    \mathbf{y} = \mathbf{W}_2 \cdot \text{ReLU}\Big( \mathbf{W}_1 \cdot \text{ReLU}(\mathbf{W}_0 \cdot \tilde{\mathbf{x}} + \mathbf{b}_1) + \mathbf{b}_2 \Big) + \mathbf{b}_3
    \label{eq:mlp_inference}
\end{equation}

โดยมีรายละเอียดเวกเตอร์อินพุต $\mathbf{x} \in \mathbb{R}^{10}$ และเวกเตอร์ผลลัพธ์ $\mathbf{y} \in \mathbb{R}^5$ ดังแสดงในตารางที่ \ref{tab:tinyml_io}

\begin{table}[htbp]
\centering
\small
\begin{tabularx}{\textwidth}{c p{4.2cm} Y c}
\toprule
\textbf{มิติ} & \textbf{ตัวแปรขาเข้า (Inputs: $\mathbf{x}$)} & \textbf{ตัวแปรเป้าหมายที่ชดเชย (Outputs: $\mathbf{y}$)} & \textbf{หน่วย} \\
\midrule
1 & Raw Nitrogen (จากรีจิสเตอร์ Modbus) & True Available Nitrogen ($N$) & $\text{mg/kg}$ \\
2 & Raw Phosphorus (จากรีจิสเตอร์ Modbus) & True Available Phosphorus ($P$) & $\text{mg/kg}$ \\
3 & Raw Potassium (จากรีจิสเตอร์ Modbus) & True Exchangeable Potassium ($K$) & $\text{mg/kg}$ \\
4 & Raw Electrical Conductivity ($\text{EC}$) & True Soil pH (ชดเชยอุณหภูมิ Nernst) & $\text{pH}$ \\
5 & Raw Soil Probe Moisture ($\%$) & True Volumetric Moisture (ผสาน 2 เซนเซอร์) & $\%$ \\
6 & Soil Temperature ($T_{\text{soil}}$ จากหัววัด) & -- & $^\circ\text{C}$ \\
7 & Air Temperature ($T_{\text{air}}$ จาก SHT45) & -- & $^\circ\text{C}$ \\
8 & Air Relative Humidity ($\text{RH}_{\text{air}}$ จาก SHT45) & -- & $\%$ \\
9 & Raw Capacitive ADC (จาก Soil Stick) & -- & $0-4095$ \\
10 & Raw Soil pH (จากรีจิสเตอร์ Modbus) & -- & $\text{pH}$ \\
\bottomrule
\end{tabularx}
\caption{รายละเอียดตัวแปรอินพุตและเอาต์พุตของโมเดล TinyML Soil Neural Calibrator}
\label{tab:tinyml_io}
\end{table}

\subsection{ผลการฝึกสอนและประเมินความแม่นยำของโมเดล}

สคริปต์ฝึกสอน \texttt{scripts/train\_soil\_calibrator.py} จำลองชุดข้อมูลฟิสิกส์เคมีดินจำนวน 4,000 ตัวอย่าง ตามสัมประสิทธิ์การเคลื่อนที่ของไอออนและปฏิกิริยาเคมีดินจริง ทำการปรับมาตรฐานการแจกแจงข้อมูลด้วย \texttt{StandardScaler} และฝึกด้วยอัลกอริทึม Adam พร้อม L2 Regularization ผลการทดสอบบนชุดข้อมูลทดสอบ (Hold-out Test Set 20\%) ปรากฏดังตารางที่ \ref{tab:tinyml_metrics}

\begin{table}[htbp]
\centering
\small
\begin{tabularx}{\textwidth}{Y c c c}
\toprule
\textbf{คุณสมบัติดินที่ชดเชย (Target Parameter)} & \textbf{สัมประสิทธิ์การตัดสินใจ ($R^2$)} & \textbf{RMSE} & \textbf{MAE} \\
\midrule
ไนโตรเจนที่ใช้ประโยชน์ได้ (Available Nitrogen) & $0.9629$ & $13.12\,\text{mg/kg}$ & $9.49\,\text{mg/kg}$ \\
ฟอสฟอรัสที่ใช้ประโยชน์ได้ (Available Phosphorus) & $0.8087$ & $9.32\,\text{mg/kg}$ & $6.76\,\text{mg/kg}$ \\
โพแทสเซียมที่แลกเปลี่ยนได้ (Exchangeable K) & $0.9639$ & $22.32\,\text{mg/kg}$ & $16.32\,\text{mg/kg}$ \\
ความเป็นกรด-ด่างของดิน (Soil pH) & $0.9890$ & $0.12\,\text{pH}$ & $0.10\,\text{pH}$ \\
ปริมาณความชื้นดินเนื้อแท้ (Volumetric Moisture) & $0.9970$ & $1.21\,\%$ & $0.96\,\%$ \\
\bottomrule
\end{tabularx}
\caption{ดัชนีชี้วัดความแม่นยำของโมเดล TinyML Soil Neural Calibrator บนชุดทดสอบ}
\label{tab:tinyml_metrics}
\end{table}

\subsection{การแปลงเป็น C++ Inference Engine และการทำงานบน ESP32-S3}

น้ำหนักโครงข่ายประสาทเทียมทั้งหมด (549 พารามิเตอร์) ถูกแปลงเป็นไฟล์เฮดเดอร์ C++ ภาษา C++ สแตนดาร์ด \texttt{gravity/include/SoilNeuralCalibrator.h} โดยใช้เทคนิค \textbf{Zero Dynamic Memory Allocation} เพื่อป้องกันการเกิดภาวะหน่วยความจำแตกกระจาย (Heap Fragmentation) บนไมโครคอนโทรลเลอร์:

\begin{cppbox}[โครงสร้าง C++ Inference Engine ของ SoilNeuralCalibrator บน ESP32-S3]
#include <Arduino.h>
#include <math.h>

struct SoilAICalibratedData {
    float nitrogen;     // Calibrated Nitrogen (mg/kg)
    float phosphorus;   // Calibrated Phosphorus (mg/kg)
    float potassium;    // Calibrated Potassium (mg/kg)
    float ph;           // Nernst temperature-compensated pH
    float moisture;     // Dual-sensor fused volumetric moisture (%)
    float confidence;   // Inference confidence score (0.0 - 1.0)
};

class SoilNeuralCalibrator {
public:
    static SoilAICalibratedData predict(
        float rawN, float rawP, float rawK, float rawEC, float rawMoist,
        float soilTemp, float airTemp, float airRH, uint16_t stickADC, float rawPH
    ) {
        float x[10] = { rawN, rawP, rawK, rawEC, rawMoist, soilTemp, airTemp, airRH, (float)stickADC, rawPH };
        
        // 1. Feature Normalization: z = (x - mean) / std
        float x_norm[10];
        for (int i = 0; i < 10; ++i) {
            x_norm[i] = (x[i] - SoilNeuralWeights::X_MEAN[i]) / SoilNeuralWeights::X_SCALE[i];
        }

        // 2. Hidden Layer 1: Dense(10 -> 16, Activation: ReLU)
        float h1[16];
        for (int j = 0; j < 16; ++j) {
            float sum = SoilNeuralWeights::BIAS_1[j];
            for (int i = 0; i < 10; ++i) sum += x_norm[i] * SoilNeuralWeights::W0[i][j];
            h1[j] = (sum > 0.0f) ? sum : 0.0f;
        }

        // 3. Hidden Layer 2: Dense(16 -> 16, Activation: ReLU)
        float h2[16];
        for (int j = 0; j < 16; ++j) {
            float sum = SoilNeuralWeights::BIAS_2[j];
            for (int i = 0; i < 16; ++i) sum += h1[i] * SoilNeuralWeights::W1[i][j];
            h2[j] = (sum > 0.0f) ? sum : 0.0f;
        }

        // 4. Output Layer: Dense(16 -> 5, Linear) and Denormalization
        SoilAICalibratedData out;
        // out.nitrogen, out.phosphorus, out.potassium, out.ph, out.moisture
        return out;
    }
};
\end{cppbox}

\subsection{การบูรณาการระบบแสดงผลและการส่งผ่านโทรมาตร (Telemetry Pipeline)}

เฟิร์มแวร์หลักของบอร์ดเชื่อมโยงผลลัพธ์ของโมเดล TinyML เข้าสู่ทุกกระบวนการของระบบ:
\begin{itemize}
    \item \textbf{หน้าจอแดชบอร์ด LCD (DisplayManager.cpp):} การ์ดที่ 4 (Deep Soil NPK \& pH) มีการประดับป้ายเรืองแสง \texttt{[AI-Edge]} โดยตัวเลขแสดงผล pH และแคปซูลเม็ดยา 3 สี (N, P, K) จะดึงค่าที่ผ่านการคำนวณของ AI ไปแสดงผลโดยตรง
    \item \textbf{หน้าต่างรายละเอียดเขตรากพืช (PAGE\_DETAIL\_SOIL7):} แสดงผลเปรียบเทียบระหว่างค่าดิบจากฮาร์ดแวร์ (Raw Sensor) กับค่าที่ผ่านการชดเชยอุณหภูมิและความชื้นของ TinyML เพื่อความโปร่งใสในการตรวจสอบ
    \item \textbf{การส่งข้อมูลคลาวด์ (CloudDataManager.cpp):} ข้อมูล \texttt{ai\_calibrated} จะถูกบรรจุลงในโครงสร้าง JSON ส่งตรงเข้าสู่ FastAPI และ Firebase Realtime Database แบบเรียลไทม์
\end{itemize}

